\subsection{Q 上的绝对值}

命题 4. 设 \(f\) 是 \(\mathbf{Q}\) 到 \(\mathbf{R}_+\) 的一个映射，且属于 \(\mathcal{V}(\mathbf{Q})\) 。那么：

\begin{enumerate}
\def\labelenumi{(\roman{enumi})}
\item
  或者 \(f\) 是 \(\mathbf{Q}\) 上的非真绝对值。
\item
  或者存在实数 \(a\) 与素数 \(p\) ，使得 \(0 < a < 1\) 且 \(f = a^{v_p}\) ，其中 \(v_p\) 是 \(p\) -进赋值。
\item
  或者存在 \(s > 0\) ，使得 \(f(x) = |x|^s\) 对一切 \(x \in \mathbf{Q}\) 成立。
\end{enumerate}

在情形 (iii) 中，\(f\) 是 \(\mathbf{Q}\) 上的绝对值，当且仅当 \(0 < s \leqslant 1\) 。

首先设 \(f(n) \leqslant 1\) 对每个整数 \(n > 0\) 成立。由第 2 节命题 3，存在实数 \(b\) 与赋值 \(v\) （\(\mathbf{Q}\) 上的），使得 \(0 < b < 1\) 且 \(f = b^v\) 。而我们知道（§ 3 第 4 节例 4），\(\mathbf{Q}\) 上（在等价的意义下）仅有的赋值是非真赋值与诸 \(p\) -进赋值 \(v_p\) ；因此或者出现情形 (i)，或者出现情形 (ii)。

从现在起设存在整数 \(h > 0\) 使得 \(f(h) > 1\) ；由第 1 节命题 2 的推论 2，存在数 \(\rho > 0\) 使得 \(f^{\rho}\) 是绝对值；记

\[
g (x) = \rho \log (f (x)) / \log | x |
\]

对每个有理数 \(x \neq 0\) 成立。设 \(a, b\) 是两个 \(\geqslant 2\) 的整数；对每个整数 \(n \geqslant 2\) ，用 \(q(n)\) 表示 \(n\) . log \(a / \log b\) 的整数部分，换句话说即最小的整数 \(m\) ，使 \(a^n < b^{m+1}\) ；于是 \(a^n\) 以 \(b\) 为底的展开为

\[
a ^ {n} = c _ {0} + c _ {1} b + \dots + c _ {q (n)} b ^ {q (n)}\tag{2}
\]

其中 \(0 \leqslant c_{i} < b\) 对 \(0 \leqslant i \leqslant q(n)\) 成立。由于 \(f^{\rho}\) 是绝对值，\(f^{\rho}(c_{i}) \leqslant c_{i} \leqslant b\) ，于是由 (2) 推出

\[
\begin{array}{r l} (f (a)) ^ {n \rho} = (f (a ^ {n})) ^ {\rho} & \leqslant b (1 + (f (b)) ^ {\rho} + \dots + (f (b)) ^ {\rho q (n)}) \\ & \leqslant b (q (n) + 1) (\sup (1, (f (b)) ^ {\rho})) ^ {q (n)}. \end{array}
\]

对这个不等式的两边取对数并除以 \(n.\log a\) ，得到

\[
g (a) \leqslant \frac {\log b}{n . \log a} + \frac {\log (q (n) + 1)}{q (n)} \cdot \frac {q (n)}{n . \log a} + \frac {\sup (0 , \rho \log f (b))}{\log a} \cdot \frac {q (n)}{n}.\tag{3}
\]

现在注意，当 \(n\) 趋于 \(+\infty\) 时，\(q(n) / n\) 趋于 \(\log a / \log b\) ；因此 \(q(n)\) 趋于 \(+\infty\) 且

\[
\log (q (n) + 1) / q (n)
\]

趋于 0（《实变函数》第三章 § 2 第 1 节）。在 (3) 中取极限，得到

\[
g (a) \leqslant \frac {\sup (0 , \rho \log f (b))}{\log b} = \sup (0, g (b)).\tag{4}
\]

但 \(f(h) > 1\) ，于是 \(g(h) > 0\) ；若在 (4) 中把 \(a\) 换成 \(h\) ，就得到 \(\sup(0, g(b)) > 0\) ，从而

\[
\sup (0, g (b)) = g (b).
\]

于是，对任意至少等于 2 的整数 \(a\) 、\(b\) ，有 \(g(a) \leqslant g(b)\) ，从而 \(g(a) = g(b)\) （交换 \(a\) 与 \(b\) 的角色）。换言之，存在常数 \(\lambda\) 使得 \(g(a) = \lambda\) 对每个整数 \(a \geqslant 2\) 成立；若记 \(s = \lambda / \rho\) ，则 \(f(a) = |a|^s\) 对每个整数 \(a \geqslant 2\) 成立。由于 \(f(xy) = f(x)f(y)\) 且 \(f(-x) = f(x) = |x|^s\) 对一切 \(x \in \mathbf{Q}\) 。最后，若 \(0 < s \leqslant 1\) ，我们知道 \(x \mapsto |x|^s\) 是绝对值（《一般拓扑学》第九章 §3 第 2 节）；反之，若 \(s\) 使 \(x \mapsto |x|^s\) 是 \(\mathbf{Q}\) 上的绝对值，则 \((1 + 1)^s \leqslant 1^s + 1^s\) ，即 \(2^s \leqslant 2\) ，于是 \(s \leqslant 1\) 。

